% ============================================================================
% 2.8 指令解析：把空格位置变成 token 指针数组
% ============================================================================

\texttt{find\_two\_spaces} 告诉我们"第一个空格在哪、第二个空格在哪"，
但要真正执行命令，我们还需要把\textbf{各个 token 的指针保存下来}，
后面按序号取回、并做指令比对。这就是 \texttt{instr.s} 要做的事：

\begin{enumerate}
    \item 分配一块内存，里面放 4 个 8 字节指针（最多 4 个 token）
    \item 把 [0, 第一个空格)、[第一个空格+1, 第二个空格)、[第二个空格+1, 末尾)
        三段分别复制到新分配的内存中（每段末尾补 0）
    \item 把这 4 个指针的\textbf{地址}放到全局的 \texttt{instr\_addrs} 数组里，
        按序号保存，供后续读取
\end{enumerate}

\subsection{数据区：instr\_addrs 全局指针数组}

\texttt{instr.s} 开头定义了一个与 \texttt{string\_array.s} 类似的平行数组，
区别是每项是一个"指针数组的头"（也就是一块包含 4 个指针的内存）：

\begin{lstlisting}[style=asm, caption=instr.s — 数据区定义]
.include "src/memory.inc"

.data
.align 4
.extern heap
.extern heap_end
.equ MAX_INSTR, 64             ; 最多保存 64 条已解析的指令
instr_addrs: .space MAX_INSTR*8  ; 每项 8 字节：指向一条 token 数组的指针
count_instr: .word 0             ; 当前已存入多少条

.text
.extern memory_init
.extern memory_alloc
.extern memory_free

.global store_instr_addrs
.global get_instr_addrs
\end{lstlisting}

布局上，\texttt{instr\_addrs} 是一个指针数组：
第 0 项指向第 0 条指令的 4 个 token 指针块，
第 1 项指向第 1 条指令的 4 个 token 指针块，以此类推。

每条指令的"4 个 token 指针块"是这样一块内存：

\begin{center}
\begin{tabular}{|l|l|}
\hline
偏移 0  & token\_0 的指针（如 "add" 的首地址） \\
\hline
偏移 8  & token\_1 的指针（如 "file1" 的首地址） \\
\hline
偏移 16 & token\_2 的指针（如 "file2" 的首地址） \\
\hline
偏移 24 & token\_3 的指针（保留 / 未使用） \\
\hline
\end{tabular}
\end{center}

\subsection{store\_instr\_addrs：把一条命令切成 token 并存下}

它接收 4 个参数：
\begin{itemize}
    \item \texttt{x0} = 第一个空格位置
    \item \texttt{x1} = 第二个空格位置
    \item \texttt{x2} = 原字符串的起始地址
    \item \texttt{x3} = 原字符串的长度
\end{itemize}

首先检查计数器，然后分配一块 32 字节（4 × 8）的内存用于存放 4 个 token 指针，
接着分 3 次把每个 token 复制到新分配的独立内存里，
每段末尾补 0 做 C 风格字符串结尾。最后把这块 4 指针的地址存入全局的
\texttt{instr\_addrs[count\_instr]}。

\begin{lstlisting}[style=asm, caption=instr.s — store\_instr\_addrs（入栈与检查）]
store_instr_addrs:
    stp x19, x20, [sp, #-96]!
    stp x21, x22, [sp, #16]
    stp x23, x24, [sp, #32]
    stp x25, x26, [sp, #48]
    stp x27, x28, [sp, #64]
    stp x29, x30, [sp, #80]

    mov x19, x0                 ; x19 = 第一个空格位置
    mov x20, x1                 ; x20 = 第二个空格位置
    mov x21, x2                 ; x21 = 字符串起始地址
    mov x22, x3                 ; x22 = 字符串总长度

    adrp x25, count_instr@PAGE
    add  x25, x25, count_instr@PAGEOFF
    ldr  w26, [x25]             ; 读当前计数
    cmp  w26, #MAX_INSTR
    b.hs skip_store             ; 已达上限，跳过

    mov  x0, #32
    bl   memory_alloc           ; 分配 32 字节 = 4 个指针
    mov  x24, x0                ; x24 = token 指针数组的头

    adrp x23, instr_addrs@PAGE
    add  x23, x23, instr_addrs@PAGEOFF
\end{lstlisting}

接下来三段几乎对称的代码分别把 token\_0、token\_1、token\_2 复制出来：

\begin{lstlisting}[style=asm, caption=instr.s — 解析并复制 3 个 token]
; ---- token 0: [0, x19) ----
parse_instr1:
    mov x0, x19                 ; 长度 = 第一个空格位置
    add x0, x0, #1              ; +1 给结尾 0 留位置
    bl  memory_alloc            ; 分配新内存存放这一段
    mov x27, x0                 ; x27 = 目标地址
    mov x28, #0
extract_byte_loop1:
    cmp x28, x19
    b.ge extract_done1
    ldrb w29, [x21, x28]
    strb w29, [x27, x28]
    add x28, x28, #1
    b extract_byte_loop1
extract_done1:
    strb wzr, [x27, x28]        ; 末尾补 0
    str  x27, [x24]             ; token 0 指针写入 instr_addrs[0]

; ---- token 1: (x19, x20) ----
parse_instr2:
    sub x0, x20, x19            ; 长度 = 第二个空格 - 第一个空格
    add x0, x0, #1
    bl  memory_alloc
    mov x27, x0
    mov x28, #0
    add x29, x19, #1            ; 从第一个空格之后开始读
extract_byte_loop2:
    cmp x29, x20
    b.ge extract_done2
    ldrb w30, [x21, x29]
    strb w30, [x27, x28]
    add x28, x28, #1
    add x29, x29, #1
    b extract_byte_loop2
extract_done2:
    strb wzr, [x27, x28]
    str  x27, [x24, #8]         ; token 1 指针写入 instr_addrs[1]

; ---- token 2: (x20, end) ----
parse_instr3:
    sub x0, x22, x20            ; 长度 = 总长 - 第二个空格
    add x0, x0, #1
    bl  memory_alloc
    mov x27, x0
    mov x28, #0
    add x29, x20, #1            ; 从第二个空格之后开始
extract_byte_loop3:
    cmp x29, x22
    b.ge extract_done3
    ldrb w30, [x21, x29]
    strb w30, [x27, x28]
    add x28, x28, #1
    add x29, x29, #1
    b extract_byte_loop3
extract_done3:
    strb wzr, [x27, x28]
    str  x27, [x24, #16]        ; token 2 指针写入 instr_addrs[2]

    str  x27, [x24, #24]        ; token 3 保留（与 token 2 同值，方便调试）
\end{lstlisting}

最后把这块 4 指针数组的头地址写入全局 \texttt{instr\_addrs}，计数器加一：

\begin{lstlisting}[style=asm, caption=instr.s — 收尾：写入全局指针数组]
    str  x24, [x23, x26, lsl #3]  ; instr_addrs[count_instr] = x24

    mov  w0, w26                   ; 返回值 = 这条指令的序号
    add  w26, w26, #1
    str  w26, [x25]                ; count_instr++

instr_addrs_done:
    ldp x29, x30, [sp, #80]
    ldp x27, x28, [sp, #64]
    ldp x25, x26, [sp, #48]
    ldp x23, x24, [sp, #32]
    ldp x21, x22, [sp, #16]
    ldp x19, x20, [sp], #96
    ret

skip_store:
    mov x0, #-1                    ; 失败返回 -1
    b instr_addrs_done
\end{lstlisting}

\subsection{get\_instr\_addrs：按序号取回 token 指针数组}

保存之后，我们希望按"这是第几条指令"把 4 个 token 的指针数组读回来。
这是写入的反向操作：

\begin{lstlisting}[style=asm, caption=instr.s — get\_instr\_addrs]
get_instr_addrs:
    stp x29, x30, [sp, #-32]!
    stp x27, x28, [sp, #16]

    mov  x28, x0                   ; x28 = 要读取的指令序号
    adrp x29, instr_addrs@PAGE
    add  x29, x29, instr_addrs@PAGEOFF
    ldr  x0, [x29, x28, lsl #3]    ; x0 = instr_addrs[序号]

get_instr_addrs_done:
    ldp x27, x28, [sp, #16]
    ldp x29, x30, [sp], #32
    ret
\end{lstlisting}

调用它之后，\texttt{x0} 就是一块 32 字节的内存，
里面依次放着 token\_0、token\_1、token\_2、token\_3 的字符串指针。
主程序里的 \texttt{get\_instr\_0 / get\_instr\_1 / get\_instr\_2 / get\_instr\_3}
就是把 \texttt{x0} 当作"指针数组"来用的例子。

\subsection{为什么要单独复制 token？}

一个很自然的问题：\emph{既然原字符串已经在堆里了，
为什么不直接保存"偏移 + 长度"，而是额外分配内存把每个 token 复制一份？}

原因有两个：

\begin{enumerate}
    \item \textbf{C 风格字符串方便后续比对}。后面如果想写
        \texttt{strcmp("add", token0)} 这种比较，要求被比较的字符串
        在自己的末尾处有一个 \texttt{0} 字节。原字符串在 token 交界处
        是空格而不是 \texttt{0}，所以必须复制一份并补 \texttt{0}。
    \item \textbf{后续可能修改或释放原字符串}。一旦原字符串所在的行
        被释放或被覆盖，我们保存在 \texttt{instr\_addrs} 里的"指向原字符串
        内部的指针"就会变成野指针。复制之后 token 各自独立，不再依赖原件。
\end{enumerate}

代价是每条指令会多 3 次 \texttt{memory\_alloc}，
但我们的堆有 4096 字节，一个 token 通常只有几个字符，在实验阶段完全够用。
